% F085｜真实数学图；代表函数不构成抽象解的唯一性证明。
\begin{center}
\begin{tikzpicture}
\begin{axis}[width=12.4cm,height=5.1cm,axis lines=middle,
 xmin=-1.4,xmax=3.5,ymin=-0.5,ymax=9.5,
 xtick={-1,0,1,2,3},xticklabels=\empty,ytick={0,2,4,8},xlabel={$t$},ylabel={$f(t)$},
 title={指数曲线 $f(t)=2^t$},title style={font=\small,yshift=3pt},
 tick label style={font=\small},label style={font=\small},
 ylabel style={at={(axis description cs:0,1)},anchor=south east,xshift=-18pt},
 clip=false,grid=major,grid style={black!8},enlargelimits=false,
 every axis plot/.append style={thick},samples=161]
\addplot[CExplanation,very thick,domain=-1.2:3.15] {2^x};
\draw[dashed,black!35] (axis cs:1,0) |- (axis cs:0,2);
\draw[dashed,black!35] (axis cs:2,0) |- (axis cs:0,4);
\draw[dashed,black!35] (axis cs:3,0) |- (axis cs:0,8);
\addplot[only marks,mark=*,mark size=2pt,orange!80!black] coordinates {(1,2) (2,4) (3,8)};

\node[font=\small,anchor=north,yshift=-3pt] at (axis cs:-1,-0.5) {$-1$};
\node[font=\small,anchor=north,yshift=-3pt] at (axis cs:1,-0.5) {$1$};
\node[font=\small,anchor=north,yshift=-3pt] at (axis cs:2,-0.5) {$2$};
\node[font=\small,anchor=north,yshift=-3pt] at (axis cs:3,-0.5) {$3$};

\end{axis}
\end{tikzpicture}
\par\small 横坐标 $1+2=3$，对应高度从 $2,4$ 相乘得到 $8$；曲线经过 $(0,1)$。
\end{center}
